\chapter{Repositorio de código}
\label{appendex:repositorio}

\section{Ubicación del repositorio}
El código de este proyecto se encuentra disponible públicamente en el repositorio GitHub del alumno en el siguiente enlace:\newline

\strut\hspace{\parskip}\url{https://github.com/ThAlsay/Hephaestus}

\section{Organización del repositorio}
El código del repositorio se organiza en las siguientes carpetas:

\begin{itemize}
	\item Database: Código SQL para la creación del esquema de la base de datos y la inserción de los datos necesarios para jugar a Alfa.
	\item Docs: Página web para el desarrollo de juegos haciendo uso de Hephaestus.
	\item Document: Este documento.
	\item Manuals: Documentos de las prácticas.
	\item Stages: Código de ejemplo y soluciones para las prácticas. Todo ello escrito en Java. Contiene también la biblioteca necesaria para que se pueda ejecutar.
	\item Tfg\_umbrella: Contiene todo el código Elixir. Esto incluye al motor Hephaestus y al juego didáctico Alfa.
\end{itemize}
